\documentclass[11pt,a4paper]{article}
\usepackage[T1]{fontenc}
\usepackage{lmodern}
\usepackage[margin=27mm]{geometry}
\usepackage{amsmath,amssymb,amsthm,mathtools}
\usepackage{microtype}
\usepackage{booktabs,array}
\usepackage{enumitem}
\usepackage[hidelinks]{hyperref}
\setlength{\parskip}{5pt plus 1pt}
\setlength{\parindent}{0pt}
\setlength{\emergencystretch}{3em}
\setlist{itemsep=2pt,topsep=3pt}
\raggedbottom
\newtheorem{theorem}{Theorem}[section]
\newtheorem{lemma}[theorem]{Lemma}
\newtheorem{corollary}[theorem]{Corollary}
\newtheorem{proposition}[theorem]{Proposition}
\theoremstyle{definition}
\newtheorem{definition}[theorem]{Definition}
\newtheorem{example}[theorem]{Example}
\theoremstyle{remark}
\newtheorem{remark}[theorem]{Remark}
\newcommand{\Z}{\mathbb{Z}}
\newcommand{\N}{\mathbb{N}}
\newcommand{\Q}{\mathbb{Q}}
\newcommand{\F}{\mathbb{F}}
\newcommand{\ps}[1]{#1[[x]]}
\newcommand{\coeff}[2]{[x^{#1}]\,#2}
\newcommand{\hl}[1]{\textbf{#1}}
\newcommand{\para}[1]{\par\medskip\noindent\textbf{#1}\enspace}
\hypersetup{pdftitle={Two congruences modulo 3 for the series of OEIS A389475},pdfauthor={Dev D. Goyal},pdfsubject={OEIS A389475}}
\title{Two congruences modulo $3$ for the series with $A(3x^2)=2x^2+A^2-2A^3$\\[6pt]\large A proof of two conjectures in OEIS A389475}
\author{Dev D. Goyal}
\date{October 2026}
\begin{document}
\maketitle
\begin{abstract}
The sequence A389475 of the On-Line Encyclopedia of Integer Sequences (OEIS) lists the coefficients of the power series
$A(x)=x+x^2+7x^3+17x^4+\cdots$ that satisfies $A(3x^2)=2x^2+A(x)^2-2A(x)^3$. Paul~D.~Hanna conjectured two congruences
modulo $3$: $A(x)^3\equiv A(x^3)$ and $A(x)^2\equiv x^2+2A(x^3)$. We prove both. We also prove that the equation has
exactly one solution with integer coefficients, which is not obvious, because computing the coefficients from the
equation requires a division by $2$ at every step. The first congruence is an instance of a classical identity that
holds for every power series with integer coefficients. The second follows from the first, because modulo $3$ the left
side $A(3x^2)$ of the equation vanishes. The results have been checked by computer and proved formally in the Lean proof
assistant. Finally, by a short argument that is not part of the formal verification, we note that modulo $9$ the
equation determines $A$ completely, as the inverse of an explicit series.
\end{abstract}
\begin{center}\begin{minipage}{0.88\linewidth}\small
\textbf{How this paper was produced.} This paper reports the result of an autonomous AI research effort.
ChatGPT (OpenAI), using the GPT-6 model with Pro reasoning and working from prompts by Dev~D.~Goyal, found the proof, wrote its Lean~4 formalization, and wrote a first draft of the paper.
Dev~D.~Goyal chose the problems and asked for checks and revisions, but did not supply proof ideas.
Claude Opus~5.5, used with Claude Code, then checked the work independently: it rebuilt the formal proofs from
source and confirmed that they rely on no unproved assumptions, recomputed the sequence from its definition and tested
every claim against the published data, and read every step of the proofs. It also rewrote this paper for clarity.
Dev~D.~Goyal has verified the results and is responsible for them. The paper has not been peer reviewed.
The section ``About this paper'' at the end gives the details.
\end{minipage}\end{center}

\section{Introduction}

\subsection{The sequence}\label{sec:sequence}
The \emph{On-Line Encyclopedia of Integer Sequences} (OEIS)~\cite{oeisdb} is a large public database of integer
sequences, used by mathematicians and scientists to identify sequences and to record what is known about them. Many
entries also record \emph{conjectures}: patterns that the contributor observed in the data but could not prove.

This paper is about entry A389475~\cite{oeis}, contributed by Paul~D.~Hanna in October 2025. It defines a power series
$A(x)=a(1)\,x+a(2)\,x^2+a(3)\,x^3+\cdots$ with integer coefficients and $a(1)=1$ by the equation
\begin{equation}\label{eq:oeis}
A(3x^2)=2x^2+A(x)^2-2A(x)^3 .
\end{equation}
On the left, $A(3x^2)$ is the series $A$ with $x$ replaced by $3x^2$, so its coefficient of $x^{2n}$ is $3^n a(n)$
(Section~\ref{sec:background} explains such substitutions). On the right, $A(x)^2$ and $A(x)^3$ are ordinary powers.
Theorem~\ref{thm:exist} shows that there is exactly one such series. Its first coefficients, starting at $n=1$, are
1, 1, 7, 17, 147, 445, 2199, 9993, 61837, 281417, \dots{} The entry reads, in part:
\begin{quote}\small\ttfamily\raggedright G.f. A(x) satisfies A(3*x\^{}2) = 2*x\^{}2 + A(x)\^{}2 - 2*A(x)\^{}3.\par\smallskip (4.a) A(x)\^{}3 == A(x\^{}3) (mod 3) (conjecture).\par\smallskip (4.b) A(x)\^{}2 == x\^{}2 + 2*A(x\^{}3) (mod 3) (conjecture).\par\smallskip\hfill\textrm{--- OEIS A389475, Paul D. Hanna, October 2025}\end{quote}

Here ``G.f.''\ stands for \emph{generating function}, the series whose coefficients are the terms of the sequence.

\subsection{The question}
The two conjectures compare power series modulo $3$, that is, coefficient by coefficient after reducing every
coefficient modulo $3$. In (4.a), $A(x)^3$ is the cube of $A$ and $A(x^3)$ is $A$ with $x$ replaced by $x^3$. The two
tables compare the two sides of each conjecture. Each row lists the coefficients of $x^n$, for $n=0,\dots,31$, of the
series named on the left, reduced modulo $3$. In each table the two rows agree.
\begin{center}\small\setlength{\tabcolsep}{3.2pt}
\begin{tabular}{>{\raggedright\arraybackslash}p{6.4em}>{\centering\arraybackslash}p{1.15em}>{\centering\arraybackslash}p{1.15em}>{\centering\arraybackslash}p{1.15em}>{\centering\arraybackslash}p{1.15em}>{\centering\arraybackslash}p{1.15em}>{\centering\arraybackslash}p{1.15em}>{\centering\arraybackslash}p{1.15em}>{\centering\arraybackslash}p{1.15em}>{\centering\arraybackslash}p{1.15em}>{\centering\arraybackslash}p{1.15em}>{\centering\arraybackslash}p{1.15em}>{\centering\arraybackslash}p{1.15em}>{\centering\arraybackslash}p{1.15em}>{\centering\arraybackslash}p{1.15em}>{\centering\arraybackslash}p{1.15em}>{\centering\arraybackslash}p{1.15em}}
\toprule
$n$ & 0 & 1 & 2 & 3 & 4 & 5 & 6 & 7 & 8 & 9 & 10 & 11 & 12 & 13 & 14 & 15 \\ \midrule
$A(x)^3$ & 0 & 0 & 0 & 1 & 0 & 0 & 1 & 0 & 0 & 1 & 0 & 0 & 2 & 0 & 0 & 0 \\
$A(x^3)$ & 0 & 0 & 0 & 1 & 0 & 0 & 1 & 0 & 0 & 1 & 0 & 0 & 2 & 0 & 0 & 0 \\
\bottomrule
\end{tabular}
\\[4pt]
\begin{tabular}{>{\raggedright\arraybackslash}p{6.4em}>{\centering\arraybackslash}p{1.15em}>{\centering\arraybackslash}p{1.15em}>{\centering\arraybackslash}p{1.15em}>{\centering\arraybackslash}p{1.15em}>{\centering\arraybackslash}p{1.15em}>{\centering\arraybackslash}p{1.15em}>{\centering\arraybackslash}p{1.15em}>{\centering\arraybackslash}p{1.15em}>{\centering\arraybackslash}p{1.15em}>{\centering\arraybackslash}p{1.15em}>{\centering\arraybackslash}p{1.15em}>{\centering\arraybackslash}p{1.15em}>{\centering\arraybackslash}p{1.15em}>{\centering\arraybackslash}p{1.15em}>{\centering\arraybackslash}p{1.15em}>{\centering\arraybackslash}p{1.15em}}
\toprule
$n$ & 16 & 17 & 18 & 19 & 20 & 21 & 22 & 23 & 24 & 25 & 26 & 27 & 28 & 29 & 30 & 31 \\ \midrule
$A(x)^3$ & 0 & 0 & 1 & 0 & 0 & 0 & 0 & 0 & 0 & 0 & 0 & 1 & 0 & 0 & 2 & 0 \\
$A(x^3)$ & 0 & 0 & 1 & 0 & 0 & 0 & 0 & 0 & 0 & 0 & 0 & 1 & 0 & 0 & 2 & 0 \\
\bottomrule
\end{tabular}
\end{center}
\begin{center}\small\setlength{\tabcolsep}{3.2pt}
\begin{tabular}{>{\raggedright\arraybackslash}p{6.4em}>{\centering\arraybackslash}p{1.15em}>{\centering\arraybackslash}p{1.15em}>{\centering\arraybackslash}p{1.15em}>{\centering\arraybackslash}p{1.15em}>{\centering\arraybackslash}p{1.15em}>{\centering\arraybackslash}p{1.15em}>{\centering\arraybackslash}p{1.15em}>{\centering\arraybackslash}p{1.15em}>{\centering\arraybackslash}p{1.15em}>{\centering\arraybackslash}p{1.15em}>{\centering\arraybackslash}p{1.15em}>{\centering\arraybackslash}p{1.15em}>{\centering\arraybackslash}p{1.15em}>{\centering\arraybackslash}p{1.15em}>{\centering\arraybackslash}p{1.15em}>{\centering\arraybackslash}p{1.15em}}
\toprule
$n$ & 0 & 1 & 2 & 3 & 4 & 5 & 6 & 7 & 8 & 9 & 10 & 11 & 12 & 13 & 14 & 15 \\ \midrule
$A(x)^2$ & 0 & 0 & 1 & 2 & 0 & 0 & 2 & 0 & 0 & 2 & 0 & 0 & 1 & 0 & 0 & 0 \\
$x^2+2A(x^3)$ & 0 & 0 & 1 & 2 & 0 & 0 & 2 & 0 & 0 & 2 & 0 & 0 & 1 & 0 & 0 & 0 \\
\bottomrule
\end{tabular}
\\[4pt]
\begin{tabular}{>{\raggedright\arraybackslash}p{6.4em}>{\centering\arraybackslash}p{1.15em}>{\centering\arraybackslash}p{1.15em}>{\centering\arraybackslash}p{1.15em}>{\centering\arraybackslash}p{1.15em}>{\centering\arraybackslash}p{1.15em}>{\centering\arraybackslash}p{1.15em}>{\centering\arraybackslash}p{1.15em}>{\centering\arraybackslash}p{1.15em}>{\centering\arraybackslash}p{1.15em}>{\centering\arraybackslash}p{1.15em}>{\centering\arraybackslash}p{1.15em}>{\centering\arraybackslash}p{1.15em}>{\centering\arraybackslash}p{1.15em}>{\centering\arraybackslash}p{1.15em}>{\centering\arraybackslash}p{1.15em}>{\centering\arraybackslash}p{1.15em}}
\toprule
$n$ & 16 & 17 & 18 & 19 & 20 & 21 & 22 & 23 & 24 & 25 & 26 & 27 & 28 & 29 & 30 & 31 \\ \midrule
$A(x)^2$ & 0 & 0 & 2 & 0 & 0 & 0 & 0 & 0 & 0 & 0 & 0 & 2 & 0 & 0 & 1 & 0 \\
$x^2+2A(x^3)$ & 0 & 0 & 2 & 0 & 0 & 0 & 0 & 0 & 0 & 0 & 0 & 2 & 0 & 0 & 1 & 0 \\
\bottomrule
\end{tabular}
\end{center}

\subsection{What this paper proves}
\textbf{Main result} (Theorem~\ref{thm:three}). \emph{Both conjectures are true:}
\[
A(x)^3\equiv A(x^3)\pmod 3\qquad\text{and}\qquad A(x)^2\equiv x^2+2A(x^3)\pmod 3 .
\]
We also prove that \eqref{eq:oeis} has exactly one normalized solution with integer coefficients
(Theorem~\ref{thm:exist}). This needs an argument: when the coefficients are computed from \eqref{eq:oeis} one at a
time, each new coefficient $a(n)$ first appears in degree $n+1$, in $A(x)^2$, with the factor $2$, so each step
involves a division by $2$.

The proofs are short. The first congruence holds for every power series with integer coefficients, by a classical
identity. For the second, note that every coefficient of $A(3x^2)$ is divisible by $3$, so modulo $3$ the equation says
$0\equiv2x^2+A(x)^2-2A(x)^3$; rearranging and using the first congruence gives the second.

\subsection{What was known}
The entry records the defining equation, some rewritings of it, and the two statements as conjectures. Hanna also
contributed the companion entry A389474~\cite{oeisother}, whose inverse series satisfies a closely related equation; it
is treated in a companion paper with the same method. Statement (4.a) is a direct consequence of a classical identity
(Lemma~\ref{lem:frobenius}); we include its proof for completeness and claim nothing new about it. We are not aware of
an earlier proof of (4.b) or of the existence of an integer solution.

\subsection{Why it is interesting}
Equations in which an unknown series is substituted into a simple series, and compared with powers of itself,
are a natural source of integer sequences, and Hanna has contributed many of them to the OEIS. Two questions come up
again and again for such equations: are the solutions integers, and do their remainders follow a pattern? This entry
shows both in a clean form. The existence of an integer solution depends on a division by $2$ that turns out to be
always exact, and the congruences come from the simple observation that the left side of the equation disappears
modulo $3$.

\subsection{How to read this paper}
Section~\ref{sec:background} explains the notions we need: power series, substitution, compositional inverses, and
remainders. Section~\ref{sec:results} states the results, and Section~\ref{sec:proofs} proves them.
Section~\ref{sec:checks} describes how the results were checked, by computer and by a formal proof in Lean.
Section~\ref{sec:discussion} discusses the scope of the results, some stronger statements, and open questions.
\section{Background}\label{sec:background}
This section collects the definitions and facts used later. Readers who know formal power series and congruences can skip it and return when needed.
\subsection{Formal power series}
A \emph{formal power series} is an expression
\[
U(x)=u_0+u_1x+u_2x^2+u_3x^3+\cdots ,
\]
determined by its sequence of coefficients $u_0,u_1,u_2,\dots$. No number is ever substituted for $x$, so questions
of convergence do not arise. We write $\coeff{n}{U}=u_n$ for the coefficient of $x^n$ in $U$; two series are equal
when all their coefficients are equal.

Series are added coefficient by coefficient and multiplied like polynomials:
\[
\coeff{n}{(UV)}=u_0v_n+u_1v_{n-1}+\cdots+u_nv_0 .
\]
Each coefficient of a product is a finite sum. For example, $(1-x)(1+x+x^2+x^3+\cdots)=1$, so $1+x+x^2+\cdots$ is
the \emph{multiplicative inverse} of $1-x$, written $\frac{1}{1-x}$. More generally, every series $D$ with constant term
$1$ has a unique multiplicative inverse $\frac1D$, whose coefficients are determined one at a time by $D\cdot\frac1D=1$
using only addition, subtraction and multiplication. Every fraction in this paper has a denominator of this kind, so it
denotes a power series.

We write $U=O(x^d)$ when all coefficients of $U$ below degree $d$ are zero. So
$U-V=O(x^d)$ means that $U$ and $V$ agree up to degree $d-1$. A series of the form $x+u_2x^2+u_3x^3+\cdots$, with
constant term $0$ and coefficient of $x$ equal to $1$, is called \emph{normalized}.

\subsection{Substituting one series into another}
If $V(x)$ has constant term $0$, we can substitute $V$ for $x$ in $U$:
\[
U(V(x))=u_0+u_1V(x)+u_2V(x)^2+u_3V(x)^3+\cdots .
\]
This is called \emph{composition}, also written $U\circ V$. It makes sense because $V(x)^k$ has no terms below $x^k$,
so only the finitely many terms with $k\le n$ affect the coefficient of $x^n$. For example, if $U=x+x^2$ and $V=x^2$, then
$U(V(x))=x^2+x^4$. Composition is associative: $U\circ(V\circ W)=(U\circ V)\circ W$.

Substituting a series into itself repeatedly gives its \emph{iterates}: $U^{\circ2}(x)=U(U(x))$,
$U^{\circ3}(x)=U(U(U(x)))$, and so on, with $U^{\circ0}(x)=x$ and $U^{\circ1}=U$. They must not be confused with
\emph{powers} $U^2=U\cdot U$, $U^3=U\cdot U\cdot U$. (The OEIS sometimes writes $A^k(x)$ for the $k$-th iterate. In
this paper the small circle always means an iterate, and a plain exponent always means a power.)

\subsection{Undoing a substitution}
A normalized series can always be ``undone'' by another normalized series.

\begin{lemma}[Compositional inverse]\label{lem:inverse}
For every normalized series $U$ there is exactly one normalized series $V$ with $U(V(x))=x$, and this $V$ also
satisfies $V(U(x))=x$. If $U$ has integer coefficients, so does $V$.
\end{lemma}

We call $V$ the \emph{compositional inverse} of $U$. The OEIS calls it the \emph{series reversion} of $U$. For example,
the compositional inverse of $x-x^2$ is $C(x)=x+x^2+2x^3+5x^4+14x^5+\cdots$, whose coefficients are the Catalan numbers.
Indeed $C-C^2=x$, which says exactly that substituting $C$ into $x-x^2$ gives $x$.

\begin{proof}
Write $V=x+v_2x^2+v_3x^3+\cdots$ with unknown coefficients. In
$U(V)=V+u_2V^2+u_3V^3+\cdots$, the coefficient of $x^n$ (for $n\ge2$) equals $v_n$ plus terms that involve only
$v_2,\dots,v_{n-1}$. Setting it equal to $0$ therefore determines $v_n$ from the earlier coefficients, using only
additions, subtractions and multiplications. So there is exactly one $V$ with $U\circ V=x$, and it has integer
coefficients if $U$ does. Applying the same construction to $V$ gives a normalized $W$ with $V\circ W=x$. Then
$U=U\circ(V\circ W)=(U\circ V)\circ W=W$, so $V\circ U=V\circ W=x$ as well.
\end{proof}

\subsection{Remainders and congruences}
For an integer $m\ge2$ we write $a\equiv b\pmod m$ (``$a$ is congruent to $b$ modulo $m$'') when $a-b$ is divisible by
$m$, that is, when $a$ and $b$ leave the same remainder after division by $m$. For example, $17\equiv5\pmod 4$, since both leave remainder $1$, and $-3\equiv1\pmod4$. The
remainder of $a$ on division by $m$ is called its \emph{residue} modulo $m$; it is one of $0,1,\dots,m-1$. Saying that
$a$ is odd is the same as saying $a\equiv1\pmod2$.

Addition and multiplication are compatible with congruences, so the residues $0,1,\dots,m-1$, added and multiplied
modulo $m$, form a ring, written $\Z/m\Z$. When $m=p$ is a prime we also write $\F_p$. In $\F_2=\{0,1\}$ we have $1+1=0$,
so addition and subtraction coincide.

For power series with integer coefficients, $U\equiv V\pmod m$ means that \emph{every} coefficient of $U-V$ is
divisible by $m$. Equivalently, the two series become equal when every coefficient is replaced by its residue; we call
this \emph{reducing} the series modulo $m$. The key fact is the following.

\emph{Reduction modulo $m$ is compatible with all the operations above.} To reduce a sum, a product, a substitution, or
an inverse, we may reduce the ingredients first and then compute in $\Z/m\Z$. This holds because every
coefficient of the result is obtained from finitely many coefficients of the ingredients by adding, subtracting and
multiplying, and these operations are compatible with remainders.

\para{Rings.} Several lemmas below hold over every \emph{commutative ring}, that is, every set with an addition,
subtraction and multiplication satisfying the usual rules of arithmetic (for example $\Z$, $\Q$ and $\Z/m\Z$). No
division is used; this matters because in $\Z/m\Z$ division is not always possible (for instance, $2$ has no inverse in
$\Z/4\Z$).

\subsection{Squaring modulo a prime}
\begin{lemma}[Frobenius identity]\label{lem:frobenius}
Let $p$ be a prime. For every power series $U$ with coefficients in $\F_p$, we have $U(x)^p=U(x^p)$.
\end{lemma}

\begin{proof}
In $\F_p$, the binomial theorem simplifies to $(a+b)^p=a^p+b^p$, because every binomial coefficient
$\binom pi$ with $0<i<p$ is divisible by $p$. Also $c^p=c$ for every $c\in\F_p$ (Fermat's little theorem). So
$\bigl(\sum_nu_nx^n\bigr)^p=\sum_nu_n^px^{pn}=\sum_nu_nx^{pn}=U(x^p)$. Only finitely many terms matter for any given
coefficient, so this is a polynomial identity.
\end{proof}

For $p=2$ this says that squaring a series modulo $2$ just doubles every exponent: for instance,
$(x+x^3)^2\equiv x^2+x^6\pmod2$. In particular, the square of any series has, modulo $2$, no terms of odd degree.

\section{Results}\label{sec:results}

\begin{theorem}[Exactly one integer solution]\label{thm:exist}
Equation \eqref{eq:oeis} has exactly one normalized solution with integer coefficients: the series $R_{-1}$ of
Lemma~\ref{lem:family} below.
\end{theorem}

\begin{theorem}[Hanna's conjectures]\label{thm:three}
Modulo $3$,
\[
\text{(a)}\quad A(x)^3\equiv A(x^3)\qquad\text{and}\qquad\text{(b)}\quad A(x)^2\equiv x^2+2A(x^3).
\]
\end{theorem}
Statement (a) holds for every power series with integer coefficients in place of $A$.

\section{Proofs}\label{sec:proofs}

\para{The idea in brief.} Equation \eqref{eq:oeis} is the case $s=-1$ of the equation $R(3x^2)=2x^2+R^2+2sR^3$. Lemma~\ref{lem:family}
shows that this equation has exactly one integer solution, by checking that the division by $2$ needed at each step is
always exact. For the congruences, we reduce modulo $3$: then $A(3x^2)$ vanishes, cubing becomes the substitution
$x\mapsto x^3$, and the equation can be rearranged into (b).

\subsection{An equation with an exact halving step}
Both this paper and its companion paper on OEIS A389474 rest on the following lemma, which concerns the
equation
\begin{equation}\label{eq:family}
R(3x^2)=2x^2+R(x)^2+2s\,R(x)^3 ,
\end{equation}
where $s$ is a fixed integer. Here $R(3x^2)$ is $R$ with $x$ replaced by $3x^2$, so its coefficient of $x^{2k}$ is
$3^kr_k$, and $R(x)^2$, $R(x)^3$ are ordinary powers.

\emph{Why the lemma is not obvious.} Suppose we try to compute the coefficients $r_2,r_3,\dots$ of a normalized
solution $R=x+r_2x^2+r_3x^3+\cdots$ one at a time. The coefficient $r_n$ first appears in degree $n+1$, in the square
$R(x)^2=x^2+\cdots+2r_nx^{n+1}+\cdots$, with the factor $2$. (In $R(x)^3$ it appears only from degree $n+2$ on, and in
$R(3x^2)$ only in degree $2n$.) So each step requires a division by $2$, and it is not clear that the result is an
integer. The proof shows that the number to be divided is always even.

\begin{lemma}[Exact halving]\label{lem:family}
For every integer $s$ there is exactly one normalized series $R_s$ with integer coefficients that satisfies
\eqref{eq:family}.
\end{lemma}

\begin{proof}
For a normalized series $U$ with integer coefficients, define its \emph{error}
\[
E(U)=U(x)^2+2s\,U(x)^3-U(3x^2)+2x^2 ,
\]
so that $U$ solves \eqref{eq:family} exactly when $E(U)=0$.

\emph{Step 1: the error is always even.} Reduce modulo $2$. By Lemma~\ref{lem:frobenius} with $p=2$,
$U(x)^2\equiv U(x^2)$. Since $3\equiv1$, also $U(3x^2)\equiv U(x^2)$. The terms $2s\,U^3$ and $2x^2$ vanish modulo $2$.
So $E(U)\equiv0\pmod2$: every coefficient of $E(U)$ is even. Moreover $E(U)$ has no terms of degree $0$, $1$ or $2$:
in degree $2$ its coefficient is $1-3+2=0$.

\emph{Step 2: changing one coefficient.} Let $V$ be another normalized integer series with
$V-U=c\,x^r+O(x^{r+1})$, where $r\ge2$. Then
\begin{itemize}
\item $V^2-U^2=(V-U)(V+U)=2c\,x^{r+1}+O(x^{r+2})$, because $V+U=2x+O(x^2)$;
\item $V^3-U^3=(V-U)(V^2+VU+U^2)=O(x^{r+2})$, because $V^2+VU+U^2=3x^2+O(x^3)$;
\item $V(3x^2)-U(3x^2)=O(x^{2r})$, which is $O(x^{r+2})$ because $r\ge2$.
\end{itemize}
Hence $E(V)-E(U)=2c\,x^{r+1}+O(x^{r+2})$. The same computation shows that if $V-U=O(x^{r+1})$, then
$E(V)-E(U)=O(x^{r+2})$.

\emph{Step 3: construction.} Start with $U=x$, and suppose that $E(U)=O(x^{r+1})$ for some $r\ge2$ (this holds for
$r=2$ by Step~1). Let $e$ be the coefficient of $x^{r+1}$ in $E(U)$. By Step~1, $e$ is even, so $e/2$ is an integer.
Replace $U$ by $U-(e/2)\,x^r$. By Step~2 the new error is $E(U)-e\,x^{r+1}+O(x^{r+2})=O(x^{r+2})$, and the coefficients
of $U$ below degree $r$ have not changed. Repeating this for $r=2,3,4,\dots$ fixes the coefficients of $U$ one at a
time. The resulting series $R_s$ has integer coefficients, and its error vanishes: for each $r$ it agrees with the
$r$-th approximation up to degree $r$, so by Step~2 its error agrees with that approximation's error, which is
$O(x^{r+2})$, up to degree $r+1$.

\emph{Step 4: uniqueness.} Suppose two different normalized solutions first differ in degree $r$, by $c\ne0$. Then
$r\ge2$, and by Step~2 their errors differ by $2c\ne0$ in degree $r+1$. But both errors are $0$, a contradiction.
\end{proof}

\subsection{Proofs of the results}

\begin{proof}[Proof of Theorem~\ref{thm:exist}]
Equation \eqref{eq:oeis} is equation \eqref{eq:family} with $s=-1$. Apply Lemma~\ref{lem:family}.
\end{proof}

\begin{proof}[Proof of Theorem~\ref{thm:three}]
(a) This is Lemma~\ref{lem:frobenius} with $p=3$; it holds for every series with integer coefficients.

(b) The coefficients of $A(3x^2)=\sum_na(n)\,3^nx^{2n}$ are all divisible by $3$. So, modulo $3$, equation
\eqref{eq:oeis} says $0\equiv2x^2+A(x)^2-2A(x)^3$, that is,
\[
A(x)^2\equiv-2x^2+2A(x)^3\equiv x^2+2A(x)^3\pmod 3 ,
\]
since $-2\equiv1$. By (a), $A(x)^3\equiv A(x^3)$, which gives (b).
\end{proof}

\section{Checking the results}\label{sec:checks}

A mathematical proof can contain mistakes, so the results were checked in two independent ways.

\para{Numerical check.}
We solved \eqref{eq:oeis} coefficient by coefficient with exact rational arithmetic, up to degree $89$. All
coefficients came out as integers, and they reproduce all 23 terms printed in the OEIS entry. Both congruences of
Theorem~\ref{thm:three} hold for every coefficient of degree below $90$.
The computation was done with the PARI/GP computer algebra system, using code that was written separately from the
formal proofs (file \texttt{verification/independent\_checks.gp} in the accompanying repository~\cite{repo}).
A computation can only test finitely many cases, so it is a check on the definitions and on the indexing, not a proof.

\para{Formal proof in Lean.}
Lean~\cite{lean} is a \emph{proof assistant}: a program in which mathematical statements and their proofs are written
in a precise formal language, and every step is checked mechanically down to the basic rules of logic. All results of
this paper are proved in Lean~4 (version 4.19.0) on top of Mathlib~\cite{mathlib}, Lean's library of formalized
mathematics. The statements are collected in the file \texttt{lean/Entries/A389475.lean} of the repository,
in the namespace \texttt{OEIS.A389475}. Lemma~\ref{lem:family} is \texttt{integral\_family\_existence}. The constructions are in \texttt{lean/Families/CubicReversion.lean}, and the congruences use \texttt{lean/Families/CubicFrobenius.lean} and \texttt{lean/Families/CubicConsequences.lean}.
The table lists where each result of this paper is proved.
\begin{center}\small\renewcommand{\arraystretch}{1.25}
\begin{tabular}{>{\raggedright\arraybackslash}p{0.52\linewidth}>{\raggedright\arraybackslash}p{0.4\linewidth}}
\toprule
Statement in this paper & Lean theorem(s)\\
\midrule
Lemma~\ref{lem:family}, for every integer $s$ & \texttt{integral\_\allowbreak{}family\_\allowbreak{}existence} \\
Theorem~\ref{thm:exist}: $A$ is normalized, solves \eqref{eq:oeis}, and is the only normalized integer solution & \texttt{normalized}, \texttt{source\_\allowbreak{}equation}, \texttt{source\_\allowbreak{}unique}, \texttt{source\_\allowbreak{}existence\_\allowbreak{}uniqueness} \\
Theorem~\ref{thm:three}(a), as a series and coefficient by coefficient & \texttt{cube\_\allowbreak{}mod\_\allowbreak{}three}, \texttt{cube\_\allowbreak{}coefficients} \\
Theorem~\ref{thm:three}(b), as a series and coefficient by coefficient & \texttt{square\_\allowbreak{}mod\_\allowbreak{}three}, \texttt{square\_\allowbreak{}coefficients} \\
\bottomrule
\end{tabular}
\end{center}
The Lean project was rebuilt from source for this paper, and Lean was asked to list the basic assumptions (axioms)
that each of these theorems ultimately depends on. Every one depends only on the three standard axioms of Lean's logic
(\texttt{propext}, \texttt{Classical.choice}, \texttt{Quot.sound}); no step is left unproved. A formal proof certifies
exactly the statement that Lean checks. Whether that statement says what the OEIS entry means is a separate question,
so the Lean definitions were also compared, side by side, with the definition in the OEIS entry
(Section~\ref{sec:sequence}).

\section{Discussion}\label{sec:discussion}

\para{Scope.} Theorem~\ref{thm:three} settles both conjectures of the entry completely. Statement (a) is not special to
this sequence.

\para{Modulo $9$ the series is determined.} The same observation gives more modulo $9$. Since $a(1)=1$ and $9$ divides
$3^n$ for $n\ge2$, we have $A(3x^2)\equiv3x^2\pmod 9$, and \eqref{eq:oeis} becomes
\[
A(x)^2-2A(x)^3\equiv x^2\pmod 9 .
\]
(Modulo $27$ one gets in the same way $A(x)^2-2A(x)^3\equiv x^2+9x^4$, because $a(2)=1$.) This determines $A$ modulo
$9$: \emph{$A$ is congruent modulo $9$ to the compositional inverse of $V(x)=x\,(1-2x)^{1/2}$.} Here
$(1-2x)^{1/2}=\sum_n\binom{1/2}{n}(-2x)^n$ is the binomial series; its coefficients are fractions whose denominators are
powers of $2$, and such fractions make sense modulo $9$ because $2$ has an inverse there ($2\cdot5\equiv1$).
\emph{Proof.} Modulo $9$, $V(x)^2=x^2-2x^3$, so $V(A(x))^2=A(x)^2-2A(x)^3=x^2$. The series $W=V(A(x))/x$ therefore has
constant term $1$ and $W^2=1$. Then $(W-1)(W+1)=0$, and $W+1$ has the invertible constant term $2$, so $W=1$. Thus
$V(A(x))=x$, and $A$ is the inverse of $V$ (Lemma~\ref{lem:inverse}). We checked this by computer for every coefficient
of degree below $90$. It is not part of the Lean formalization.

\para{The family.} Equation \eqref{eq:family} has an integer solution $R_s$ for every integer $s$. This entry is the case
$s=-1$; the compositional inverse of the series of the companion entry A389474 is the case $s=1$. For that entry, the
companion paper shows that $A(x)^2\equiv x^2+2x^3\pmod 9$. It would be interesting to understand the other members of
the family, and the remainders of the coefficients $a(n)$ themselves modulo higher powers of $3$.

\section*{About this paper}\label{sec:about}
\para{How the work was done.}
This paper is one of 26 papers on conjectures from the OEIS that came out of an autonomous AI research effort directed
by Dev~D.~Goyal~\cite{repo}. The roles were as follows.
\begin{itemize}
\item \emph{Dev~D.~Goyal} started the project, chose which problems to pursue, asked for proofs, formal verification and
revisions, and decided what to publish. Dev~D.~Goyal did not supply proof strategies or intermediate arguments.
\item \emph{ChatGPT} (OpenAI), using the GPT-6 model with Pro reasoning, worked autonomously from those prompts. It
searched the OEIS for suitable conjectures, found the proofs, wrote the Lean~4 formalization, and wrote a first draft
of this paper.
\item \emph{Claude Opus~5.5}, used with Claude Code, checked the work independently and then rewrote this paper. It
rebuilt the Lean project from source and confirmed that the theorems listed in Section~\ref{sec:checks} rely only on
Lean's standard axioms; compared the Lean definitions with the OEIS definition; recomputed the sequence with separately
written PARI/GP code, compared it with the terms published in the OEIS, and tested every claim of the paper on many
terms; and read every step of the proofs. It then rewrote the text to make it self-contained and readable. The
theorems and the proof strategies are ChatGPT's. In rewriting, Claude Opus~5.5 expanded the explanations and, in
places, wrote out an argument in more detail or by a more elementary route. The Discussion section was added by Claude Opus~5.5;
the further observations in it were checked by computer and are not part of the Lean formalization.
\end{itemize}
Dev~D.~Goyal has verified the results and is responsible for this paper.

\para{Review status.} This paper has not been peer reviewed, and it has not been submitted to a journal or to the OEIS.

\para{Priority.} An OEIS entry that labels a statement as a conjecture shows only that its contributor did not know
a proof; it does not show that no proof existed. During the research only limited searches for earlier work were
made. No claim of priority is made.

\para{Corrections.} Please report any error in the issue tracker of the repository~\cite{repo}, where the Lean sources and the numerical checks are also
available.

\section*{Acknowledgements}
The sequence and the conjectures are due to Paul~D.~Hanna~\cite{oeis}, who also contributed the companion entry
A389474~\cite{oeisother}. Statement (4.a) is an instance of the classical Frobenius identity and is not claimed as new.

\begin{thebibliography}{99}\small
\raggedright
\bibitem{oeis} P.~D.~Hanna, entry A389475 (G.f.\ $A(x)$ satisfies $A(3x^2)=2x^2+A(x)^2-2A(x)^3$), \emph{The On-Line Encyclopedia of Integer Sequences}, October 2025. \url{https://oeis.org/A389475}. Text as of October 2026.
\bibitem{oeisother} P.~D.~Hanna, entry A389474 (G.f.\ $A(x)$ satisfies $A(x^2+2x^3+2A(x)^2)=3A(x)^2$), \emph{The On-Line Encyclopedia of Integer Sequences}, October 2025. \url{https://oeis.org/A389474}. Text as of October 2026.
\bibitem{oeisdb} N.~J.~A.~Sloane and The OEIS Foundation Inc., \emph{The On-Line Encyclopedia of Integer Sequences}. \url{https://oeis.org}.
\bibitem{stanley} R.~P.~Stanley, \emph{Enumerative Combinatorics}, Volume~2, Cambridge University Press, 1999. (Chapter~5 treats formal power series, composition and compositional inverses.)
\bibitem{niven} I.~Niven, Formal power series, \emph{American Mathematical Monthly} 76 (1969), 871--889.
\bibitem{lean} L.~de~Moura and S.~Ullrich, The Lean~4 theorem prover and programming language, in \emph{Automated Deduction -- CADE 28}, Lecture Notes in Computer Science 12699, Springer, 2021, 625--635.
\bibitem{mathlib} The mathlib Community, The Lean mathematical library, in \emph{Proceedings of the 9th ACM SIGPLAN International Conference on Certified Programs and Proofs (CPP 2020)}, 367--381. Version used: commit \texttt{c44e0c8ee63c}.
\bibitem{repo} D.~D.~Goyal, \emph{Proofs and corrections of OEIS conjectures, with Lean~4 formalizations}, repository, 2026. \url{https://github.com/dvgyl/oeis-a396843-proof}.
\end{thebibliography}
\end{document}
